\documentclass[11pt]{article}

\usepackage[margin=1in]{geometry}
\usepackage[T1]{fontenc}
\usepackage{lmodern}
\usepackage{microtype}
\usepackage{amsmath,amssymb}
\usepackage[dvipsnames]{xcolor}
\usepackage{enumitem}
\usepackage{needspace}
\usepackage[most]{tcolorbox}
\usepackage[hidelinks]{hyperref}

\hypersetup{
  pdftitle={Extra practice for Midterm - Solutions, MATH 327, Fall 2026},
  pdfauthor={Manish Patnaik},
  pdfsubject={Midterm 1 extra practice}
}

\definecolor{courseblue}{HTML}{1F4E79}
\definecolor{lightblue}{HTML}{EAF2F8}
\definecolor{boxborder}{HTML}{B8CBDC}
\definecolor{darkgray}{HTML}{333333}
\definecolor{lightyellow}{HTML}{FFF9DC}
\definecolor{yellowborder}{HTML}{E4D58C}

\setlength{\parindent}{0pt}
\setlength{\parskip}{0.55em}
\setlist[enumerate]{leftmargin=2em, itemsep=0.25em, topsep=0.3em}

\newtcolorbox{exercise}[1]{
  enhanced,
  colback=white,
  colframe=boxborder,
  colbacktitle=lightblue,
  coltitle=courseblue,
  fonttitle=\bfseries,
  title={Exercise #1},
  boxrule=0.6pt,
  arc=1mm,
  left=2.5mm,
  right=2.5mm,
  top=1.5mm,
  bottom=1.5mm,
  before skip=0.8em,
  after skip=0.8em
}

\newtcolorbox{solution}{
  enhanced,
  breakable,
  colback=lightyellow,
  colframe=yellowborder,
  colbacktitle=lightyellow,
  coltitle=darkgray,
  fonttitle=\bfseries,
  title={Solution},
  boxrule=0.6pt,
  arc=1mm,
  left=2.5mm,
  right=2.5mm,
  top=1.5mm,
  bottom=1.5mm,
  before skip=0.2em,
  after skip=1em
}

\newcommand{\chapterheading}[1]{%
  \par\vspace{0.5em}
  {\color{courseblue}\Large\bfseries #1}\par\vspace{0.35em}
  {\color{boxborder}\hrule height 0.6pt}\vspace{0.35em}
}

\newcommand{\sectionheading}[1]{%
  \Needspace{8\baselineskip}
  \vspace{0.25em}
  {\color{darkgray}\large\bfseries #1}\par
}

\begin{document}
\begin{center}
\colorbox{lightblue}{\begin{minipage}{0.92\textwidth}
\vspace{0.5em}
{\color{courseblue}\bfseries\Large Extra practice for Midterm}\\[0.4em]
MATH 327: Algebra 1 \hfill Fall 2026\\[0.2em]
Midterm 1 \hfill Solutions
\vspace{0.5em}
\end{minipage}}
\end{center}
The first three exercises are from Michael Artin's \emph{Algebra}, second edition,
Chapter 2, Section 8 (Cosets).
\begin{exercise}{8.5}
A finite group $G$ contains an element $x$ of order $10$ and also an element
$y$ of order $6$. What can be said about the order of $G$?
\end{exercise}
\begin{solution}
By Lagrange's theorem, the orders of $\langle x\rangle$ and
$\langle y\rangle$ divide $|G|$. Hence $10\mid |G|$ and $6\mid |G|$, so
\[
\boxed{30\mid |G|.}
\]
In particular, $|G|\geq30$. Every positive multiple of $30$ can occur:
if $G=\langle g\rangle$ is cyclic of order $30m$, then $g^{3m}$ has
order $10$ and $g^{5m}$ has order $6$.
\end{solution}
\begin{exercise}{8.7}
A group $G$ of order $22$ contains elements $x$ and $y$, where $x\ne1$
and $y$ is not a power of $x$. Prove that the subgroup generated by these
elements is the whole group $G$.
\end{exercise}
\begin{solution}
Let $H=\langle x,y\rangle$. Since $H\leq G$, Lagrange's theorem gives
$|H|\in\{1,2,11,22\}$. Because $x\ne1$, we have $|H|>1$.

If $|H|$ were $2$ or $11$, then $H$ would have prime order. The
nonidentity element $x\in H$ would therefore generate $H$: its order
divides $|H|$ and is greater than $1$, so it equals $|H|$.
But then $y\in H=\langle x\rangle$ would be a power of $x$, contrary
to the hypothesis. Thus $|H|=22$, and consequently $H=G$.
\end{solution}
\begin{exercise}{8.8}
Let $G$ be a group of order $25$. Prove that $G$ has at least one subgroup
of order $5$, and that if it contains only one subgroup of order $5$,
then it is a cyclic group.
\end{exercise}
\begin{solution}
Choose $a\ne1$ in $G$. By Lagrange's theorem, its order is either $5$
or $25$. If its order is $5$, then $\langle a\rangle$ is a subgroup
of order $5$. If its order is $25$, then $a^5$ has order $5$:
$(a^5)^5=1$, while $a^5\ne1$, so its order is $5$.
Thus $G$ always has a subgroup of order $5$.

Now suppose $H$ is the unique subgroup of order $5$. Since $|H|<|G|$,
choose $b\in G\setminus H$. In particular, $b\ne1$, so its order is
$5$ or $25$. If its order were $5$, then $\langle b\rangle$ would be
a subgroup of order $5$ and hence would equal $H$, contradicting
$b\notin H$. Therefore $b$ has order $25$, so
$G=\langle b\rangle$ is cyclic.
\end{solution}
\newpage
\chapterheading{Short-answer practice}
Give the requested answers; no proof is required. Compose permutations rightmost first.
\begin{exercise}{S1}
In the additive group $\mathbb Z/20\mathbb Z$, find the order of
$\overline{8}$ and list the elements of $\langle\overline{8}\rangle$.
\end{exercise}
\begin{solution}
The order is $5$, and
\[
\langle\overline{8}\rangle
=\{\overline{0},\overline{4},\overline{8},\overline{12},\overline{16}\}.
\]
Indeed, $k\overline{8}=\overline{0}$ exactly when $20\mid8k$, or equivalently $5\mid k$.
\end{solution}
\begin{exercise}{S2}
Let $\varphi:\mathbb Z\to\mathbb Z/12\mathbb Z$ be the homomorphism
$\varphi(n)=\overline{4n}$. Describe $\ker\varphi$ and list the elements
of $\operatorname{im}\varphi$.
\end{exercise}
\begin{solution}
We have $\ker\varphi=3\mathbb Z$, because $12\mid4n$ if and only if
$3\mid n$. The image is $\{\overline{0},\overline{4},\overline{8}\}$.
\end{solution}
\begin{exercise}{S3}
Find the positive integers $d$ and $e$ such that
\[
8\mathbb Z+12\mathbb Z=d\mathbb Z,
\qquad 8\mathbb Z\cap12\mathbb Z=e\mathbb Z.
\]
\end{exercise}
\begin{solution}
$d=4$ and $e=24$. The sum is generated by $\gcd(8,12)=4$;
the intersection consists of the common multiples of $8$ and $12$.
\end{solution}
\begin{exercise}{S4}
In the additive group $\mathbb Z/12\mathbb Z$, let
$H=\langle\overline{4}\rangle$. List the elements of the coset
$\overline{1}+H$. Is $\overline{9}+H$ the same coset?
\end{exercise}
\begin{solution}
$H=\{\overline{0},\overline{4},\overline{8}\}$, so
$\overline{1}+H=\{\overline{1},\overline{5},\overline{9}\}$.
Yes: $\overline{9}-\overline{1}=\overline{8}\in H$.
\end{solution}
\newpage
\chapterheading{Short-answer practice (continued)}
\begin{exercise}{S5}
Let $\sigma=(1\ 5\ 2\ 4)(3\ 6)\in S_6$.
Find $\sigma^{-1}$, the order of $\sigma$, and $\operatorname{sgn}(\sigma)$.
\end{exercise}
\begin{solution}
$\sigma^{-1}=(1\ 4\ 2\ 5)(3\ 6)$, and
$\operatorname{ord}(\sigma)=\operatorname{lcm}(4,2)=4$.
Both the $4$-cycle and the transposition have sign $-1$, so
$\operatorname{sgn}(\sigma)=1$.
\end{solution}
\begin{exercise}{S6}
Let $G$ be a group of order $42$ and $H\leq G$ a subgroup of order $7$.
How many left cosets of $H$ are there? Can $G$ contain an element of order $8$?
\end{exercise}
\begin{solution}
There are $[G:H]=42/7=6$ left cosets. No element can have order $8$,
since an element's order divides $42$ by Lagrange's theorem, and $8\nmid42$.
\end{solution}
\begin{exercise}{S7}
Let $\psi:G\to K$ be a homomorphism, where $|G|=24$ and
$|\ker\psi|=6$. Find $|\operatorname{im}\psi|$.
If $|K|=8$, is $\psi$ surjective?
\end{exercise}
\begin{solution}
$|\operatorname{im}\psi|=|G|/|\ker\psi|=24/6=4$.
If $|K|=8$, then $\psi$ is not surjective, since its image has only $4$ elements.
\end{solution}
\begin{exercise}{S8}
Let $H=\{1,(1\ 3)\}\leq S_3$.
List the elements of the left coset $(1\ 2)H$ and of the right coset
$H(1\ 2)$. Are these cosets equal?
\end{exercise}
\begin{solution}
Composing rightmost first gives
\[
(1\ 2)H=\{(1\ 2),(1\ 3\ 2)\},\qquad
H(1\ 2)=\{(1\ 2),(1\ 2\ 3)\}.
\]
The two cosets are not equal.
\end{solution}
\newpage
\chapterheading{True or false}
Mark each statement true or false. For review, also try to give a brief reason or a counterexample.
\textbf{T1.} Every cyclic group is abelian.
\begin{solution}
\textbf{True.} If $G=\langle a\rangle$, then $a^ra^s=a^{r+s}=a^{s+r}=a^sa^r$.
\end{solution}
\textbf{T2.} Every abelian group is cyclic.
\begin{solution}
\textbf{False.} The additive group $(\mathbb Z/2\mathbb Z)\times(\mathbb Z/2\mathbb Z)$ has four elements, but each nonidentity element has order $2$.
\end{solution}
\textbf{T3.} Every left coset of a subgroup $H\leq G$ is a subgroup of $G$.
\begin{solution}
\textbf{False.} The coset $1+2\mathbb Z$ in $(\mathbb Z,+)$ does not contain the identity $0$.
\end{solution}
\textbf{T4.} If $H$ is a finite subgroup of $G$, then every left coset of $H$ has $|H|$ elements.
\begin{solution}
\textbf{True.} For any $g\in G$, the map $h\mapsto gh$ is a bijection from $H$ to $gH$.
\end{solution}
\textbf{T5.} If $\varphi:G\to K$ is a homomorphism and $a\in G$ has finite order, then $\varphi(a)$ has the same order as $a$.
\begin{solution}
\textbf{False.} The trivial homomorphism from a group of order $2$ to $\{1\}$ sends its nonidentity element to an element of order $1$. The image order need only divide the original order.
\end{solution}
\newpage
\chapterheading{True or false (continued)}
\textbf{T6.} A group homomorphism with trivial kernel is injective.
\begin{solution}
\textbf{True.} If $\varphi(x)=\varphi(y)$, then $\varphi(y^{-1}x)=1$. A trivial kernel forces $y^{-1}x=1$, hence $x=y$.
\end{solution}
\textbf{T7.} If $a\in G$ has finite order and $g\in G$, then $gag^{-1}$ has the same order as $a$.
\begin{solution}
\textbf{True.} $(gag^{-1})^m=ga^mg^{-1}$, which equals $1$ exactly when $a^m=1$.
\end{solution}
\textbf{T8.} The product of two odd permutations is odd.
\begin{solution}
\textbf{False.} Sign is multiplicative, and $(-1)(-1)=1$, so the product is even.
\end{solution}
\textbf{T9.} The set of invertible $2\times2$ real matrices is a group under matrix multiplication.
\begin{solution}
\textbf{True.} Products and inverses remain invertible, the identity matrix belongs to the set, and matrix multiplication is associative.
\end{solution}
\textbf{T10.} $4\mathbb Z\subseteq6\mathbb Z$.
\begin{solution}
\textbf{False.} $4\in4\mathbb Z$, but $4\notin6\mathbb Z$. In general, $a\mathbb Z\subseteq b\mathbb Z$ if and only if $b\mid a$.
\end{solution}
\end{document}
